\section{灵敏度分析与稳健性}\label{sec:sens}

\subsection{可测性口径：从"一定听得到"到"可能要多测一次"}

本文采用最保守口径（源到第二点的距离 $\le1000$ m）。若允许第二次测向失败后重跑，把可测性
放宽到接收半径上界 $1500$ m，则可行域从 $0.447$ km$^2$ 扩到 $2.737$ km$^2$，最优最坏直径
从 $134.08$ m 降到 $79.99$ m，最优点外移到 $r=1500$ m、$\varphi=+39.8^\circ$（"贴到最远处"），
候选弧带变成 $r\in[1262,1504]$ m、$\theta_1\pm[29.5^\circ,45.5^\circ]$、面积 $0.0931$ km$^2$。
这说明：\textbf{代价主要在"保证"上，不在"方向"上}——放宽保证能白拿 $40\%$ 的精度，但会带来
"白跑一趟"的风险，本文不采用。

\subsection{第一轮已知距离先验}

若第一轮还能给出距离区间（例如通过信号强度或接收时间给出 $d\le1000$ m 的上限），不确定集
被压短，可行域变为 $1.219$ km$^2$，最优最坏直径直接降到 $53.31$ m（最优点
$(769,639)$ m、$r=1000$ m、$\varphi=+39.7^\circ$），候选弧带 $r\in[842,1004]$ m、
面积 $0.0415$ km$^2$。可见\textbf{距离信息是最值钱的信息}：把"最远可能源从 $1500$ m 缩到
$1000$ m"这一条，就能把最坏定位直径减少六成。

\begin{table}[H]
  \centering
  \caption{三种口径下的最优解与候选区域（其余参数不变）}
  \label{tab:sens}
  \small
  \setlength{\tabcolsep}{5pt}
  \begin{tabular}{lcccccc}
    \toprule
    \textbf{口径} & $S_2^*$ / m & $r^*$ / m & $\varphi^*$ & $J^*$ / m & \textbf{可行域 / km$^2$} & \textbf{提升} \\
    \midrule
    缺省（保证 $\le1000$ m）        & $(801,605)$   & $1004$ & $+37.1^\circ$ & $134.08$ & $0.447$ & $13.4$ 倍 \\
    可测性放宽到 $1500$ m           & $(1153,959)$  & $1500$ & $+39.8^\circ$ & $79.99$  & $2.737$ & $22.5$ 倍 \\
    已知 $d\le1000$ m（距离先验）   & $(769,639)$   & $1000$ & $+39.7^\circ$ & $53.31$  & $1.219$ & $33.8$ 倍 \\
    \bottomrule
  \end{tabular}
\end{table}

\subsection{候选区域阈值}

阈值 $\eta$ 只改变弧带的宽度而不改变最优点：$\eta=10\%$ 时面积 $0.02069$ km$^2$，
$\eta=5\%$ 时 $0.00880$ km$^2$，$\eta=2\%$ 时 $0.00289$ km$^2$（表~\ref{tab:band}）。
即"宽容度"与"可现场选择的范围"之间存在直接的、单调的取舍关系，现场可以按需要的余量挑
阈值。

\subsection{检测点位置与示向度方向}

模型与算法对 $S_1$ 的位置与示向度没有任何特殊依赖。用三组完全不同的设定复算：

\begin{itemize}
  \item $S_1=(200,-300)$ m、$\theta_1=45^\circ$：$J^*=134.10$ m；
  \item $S_1=(1700,0)$ m、$\theta_1=180^\circ$：$J^*=134.10$ m；
  \item $S_1=(-900,900)$ m、$\theta_1=225^\circ$：$J^*=79.34$ m、可行域 $0.383$ km$^2$。
\end{itemize}

前两组与缺省解的 $J^*$ 完全相同（$134.10$ m 与 $134.08$ m 的差别来自 $5$ m 认证网格的步长
自适应，属于数值分辨率），说明结论具有\textbf{平移与旋转不变性}——这符合问题的几何性质：
把整幅构型一起平移旋转，直径不变。第三组之所以更好，是因为该检测点距圆域中心 $1273$ m，
示向度指向圆外，最远的那些源（最不友好的共线构型）被 $1800$ m 圆域裁掉了，不确定集本身变短，
可行域也随之变小。这提示了一个有用的工程推论：\textbf{如果检测点靠近圆域边界、且示向度朝外，
第二个检测点的精度要求会更宽松}。

\subsection{判据选择与快筛稳定性}

§\ref{sec:res} 表~\ref{tab:three} 已说明三种准则的最优点相差不足 $20$ m。判据快筛的稳定性
同样在多组口径下检验：GDOP 与精确直径的排序相关在缺省口径为 $0.998$，在可行域被拉长的
两个放宽口径下分别为 $0.942$（可测性放宽）与 $0.946$（距离先验）。$0.94$ 意味着排序略有
损失，但仍足以把几万个候选压缩到几十个，最终结论一律由精确构造复核给出——快筛只影响
"算多久"，不影响"答案是什么"。
